\documentclass[11pt]{article}

\usepackage[margin=1in]{geometry}
\usepackage{amsmath,amssymb,amsthm}
\usepackage{booktabs}
\usepackage{array}
\usepackage{enumitem}
\usepackage{natbib}
\usepackage[hidelinks]{hyperref}

\newtheorem{proposition}{Proposition}
\newcommand{\R}{\mathbb{R}}
\newcommand{\dd}{\mathrm{d}}
\newcommand{\KDM}{\ensuremath{\mathrm{KDM}}}
\newcommand{\LEDH}{\ensuremath{\mathrm{LEDH}}}
\newcommand{\GenUT}{\ensuremath{\mathrm{GenUT}}}
\newcommand{\OT}{\ensuremath{\mathrm{OT}}}

\title{Younis Kernel Mixtures and the Score of the\\
LEDH--OT--GenUT Dual-Cap Trust-Region Filter}
\author{BayesFilter research note}
\date{9 September 2026}

\begin{document}
\maketitle

\begin{abstract}
The question in this note is deliberately narrow.  Can the
kernel-density-mixture construction of Younis and Sudderth
\citep{younis2023,younis2024} help compute the score of the executed
LEDH--OT--GenUT filter with its dual-cap trust-region correction?  The answer
cannot be obtained by calling a smoothed derivative a correction to the
existing score.  A positive-bandwidth mixture changes the finite state
measure and therefore changes the scalar being differentiated.  Younis's
importance-weighted sample gradient (IWSG) gives a principled derivative for
a mixture expectation with a proposal frozen at the evaluation point; it is
not, without an additional argument, the total derivative of the
atom-based filter program.  This note defines the distinct programs, derives
the complete mixture and Gaussian-convolution differentials, identifies the
support problem for degenerate DSGE states, and sets out a bounded experiment.
The existing LEDH--OT--\GenUT{} total derivative remains the baseline.  The
implemented observation-only KDM program passes total-derivative checks for
its own scalar, but a paired linear--Gaussian experiment does not show lower
score error.  The source audit also resolves a more basic issue: Younis's KDM
is a continuous resampling law, whereas Contract--E is already the reset of
the requested filter.  Replacing Contract--E does not implement the requested
algorithm; retaining both defines a new, combined finite program.  That
combined fixed-anchor program is now implemented as an exact all-components
reference.  A later paper-and-author-code audit found that the first version
prematurely normalized the IWSG ratios.  That normalization left the next
posterior weights unchanged but subtracted a random finite-sample term from
the next log normalizer and its derivative.  The repaired primary route now
carries the raw ratio divided by the particle count, as required by the IWSG
expectation identity; the old centered route is a separately labelled
self-normalized program.  Whether the repaired derivative estimates the model
score better has not yet been tested.  Neither route is the derivative of the
unchanged atom program.
\end{abstract}

\section{The question and the target}

The repository uses \LEDH{} particle flow, PF--PF proposal weights,
entropy-regularized \OT{}, \GenUT{} moment restoration, and a dual-cap
trust-region correction.  The spelling \LEDH{} follows the repository's
Li--Coates route \citep{li2017particle}; if a different method is intended by
\lq\lq LEDG\rq\rq, its definition must replace the notation before
implementation.

The present research question is not whether a generic particle smoother can
be attached to the project.  The question is whether a Younis-style
continuous mixture can improve the computation of the score associated with
the actual finite filter.  Nemeth, Scibior--Wood, PaRIS, and Del Moral are
therefore not competing implementation directions in this note.  Corenflos
is part of the existing differentiable-\OT{} context
\citep{corenflos2021differentiable}.  These methods may delimit what is being
computed, but the proposed experiment is a KDM experiment.

There are two score notions that must be kept separate:
\begin{align}
  \ell(\theta) &= \log p_\theta(y_{1:T}), &
  s(\theta) &= \nabla_\theta \ell(\theta), \label{eq:exact-score}\\
  L_0^N(\theta;\xi) &= \text{value of the executed finite filter}, &
  S_0^N(\theta;\xi) &= \nabla_\theta L_0^N(\theta;\xi). \label{eq:finite-score}
\end{align}
The first is the exact observed-data score.  The second is the total
derivative of the finite program for a fixed observation sequence and fixed
random stream $\xi$.  The second is the immediate target of the present code
audit.  Equality between the two is not assumed at finite particle count.

The experiments below use four active labels and one historical diagnostic
label.  \texttt{ATOM-FINITE} denotes the
canonical scalar $L_0^N$ in \eqref{eq:finite-score}.  \texttt{KDM-FINITE}
denotes a scalar in which a declared positive-bandwidth operation has replaced
an atom operation, with score $\nabla_\theta L_h^N$.
\texttt{RESKDM-IWSG-FINITE} denotes the particular combined program that
retains Contract--E and then carries raw fixed-anchor full-mixture importance
weights into the next transition.  \texttt{RESKDM-SN-FINITE} denotes the
historical variant that first self-normalizes those ratios.  \texttt{MODEL-IS}
denotes an importance-sampling estimator
whose proposal may be a KDM but whose density ratio targets the underlying
state-space factor.  A correct
\texttt{MODEL-IS} ratio preserves the model normalizer in expectation under
its support assumptions; it does not make the finite Contract--E reset equal
to \texttt{ATOM-FINITE}.  These labels are not interchangeable.  In
particular, \texttt{RESKDM-SN-FINITE} is wrong relative to a claim that it
executes the raw IWSG normalizer.

\section{The executed atom program}

Let $x_{t,i}^{-}$ and $w_{t,i}^{-}$ denote the cloud and normalized weights
immediately before the observation update.  The finite program can be written
schematically as
\begin{align}
 x_{t,i}^{-}
   &= \Phi_{t,\theta}(x_{t-1,i}^{+},\xi_t), \\
 x_{t,i}^{\mathrm{out}}
   &= F_{t,\theta}^{\mathrm{LEDH}}(x_{t,i}^{-};\xi_t), \\
 a_{t,i}
   &= \log w_{t,i}^{-}
      +\log p_\theta(x_{t,i}^{\mathrm{out}}\mid x_{t-1,i}^{+})
      +\log g_\theta(y_t\mid x_{t,i}^{\mathrm{out}})
      +\log |J_{t,i}|
      -\log q_{t,\theta}(x_{t,i}^{-}\mid x_{t-1,i}^{+},y_t), \\
 Z_t
   &= \sum_{i=1}^{N}\exp(a_{t,i}), \\
 w_{t,i}^{+}
   &= \exp(a_{t,i})/Z_t, \\
 x_t^{+}
   &= R_\theta(x_t^{\mathrm{out}},w_t^{+};\xi_t), \label{eq:atom-program}\\
 L_0^N(\theta;\xi)
   &= \sum_{t=1}^{T}\log Z_t,\qquad
 S_0^N(\theta;\xi)=\nabla_\theta L_0^N(\theta;\xi).
\end{align}
Here $x_{t,i}^{\mathrm{out}}$ is the post-flow state, $J_{t,i}$ is the
forward flow Jacobian, and $q_{t,\theta}$ is the pre-flow proposal density.
$\Phi$ includes the declared transition and UKF covariance lifecycle, while
$F^{\mathrm{LEDH}}$ is the flow map; the displayed terms make the PF--PF
correction explicit.  In a
bootstrap/no-flow special case, the proposal and transition terms cancel and
$a_{t,i}=\log w_{t,i}^{-}+\log g_\theta(y_t\mid x_{t,i}^{-})$.  The reset $R$
includes the Sinkhorn/\OT{} map, Contract--E moments, \GenUT{} corrections,
and the dual-cap trust-region map.  This form matches the canonical score
route, whose weight logits contain transition, observation, log-determinant,
and proposal terms.  No such term may be silently dropped from the value or
its tangent.

For the auxiliary sidecar it is useful to expose the same code path as a
prior-observation factorization.  Define
\begin{align}
 b_{t,i} &= \log w_{t,i}^{-}
      +\log p_\theta(x_{t,i}^{\mathrm{out}}\mid x_{t-1,i}^{+})
      +\log |J_{t,i}|
      -\log q_{t,\theta}(x_{t,i}^{-}\mid x_{t-1,i}^{+},y_t),\\
 B_t &= \log\sum_i\exp(b_{t,i}),\qquad
 \bar w_{t,i}=\exp(b_{t,i}-B_t),\\
 a_{t,i} &= b_{t,i}+\log g_\theta(y_t\mid x_{t,i}^{\mathrm{out}}),\qquad
 Z_t=\sum_i\exp(a_{t,i}).
 \label{eq:prior-observation-factorization}
\end{align}
The implementation records $\bar w_t$ as
\texttt{prior\_observation\_weights}; these are the weights that may be
used with a new observation functional.  The posterior weights in
\eqref{eq:atom-program} are $w_t^+$ and must not be reused as $\bar w_t$.

The derivative in \eqref{eq:finite-score} is a total derivative.  Thus a
tangent through $x_t^{-}$, $w_t^{-}$, the observation factor, the \OT{} map,
the Contract--E residual design, the higher moments, and the dual caps is
included whenever that dependence is part of the declared finite scalar.
Holding a cloud fixed and differentiating only the observation factor is a
different partial derivative.

The executable analytical endpoint evaluates one arbitrary parameter
direction per call.  For a scalar parameter this number is the score.  For
$p>1$, a full score vector requires $p$ calls with the model tangent set to
the coordinate directions (or another spanning basis); each call is still a
total derivative through the entire finite recurrence.  The present Phase~4A
and Phase~4B LGSSM fixtures have one parameter.  No result below establishes a
multi-parameter Phase~4B assembly call site.

The single-cloud analytical score authority inspected for this note is:
\begin{itemize}[leftmargin=2em]
  \item \texttt{canonical\_value\_and\_analytical\_score}\\
        \path{bayesfilter/highdim/ledh_canonical_score_tf.py};
  \item the batch wrapper
        \texttt{canonical\_batch\_value\_score}\\
        \path{bayesfilter/highdim/ledh_canonical_batch_}\\
        \path{tf.py}.\\
        It accepts both the historical \texttt{substeps} keyword and the
        authority-aligned \texttt{flow\_substeps} alias, and rejects
        conflicting values.
\end{itemize}
The separate lower-level GenUT candidate/diagnostic is:
\begin{itemize}[leftmargin=2em]
  \item \texttt{finite\_value\_score}\\
        \path{bayesfilter/highdim/cubature_genut_filter.py}.
\end{itemize}
Calling it is not evidence that the LEDH PF--PF flow and Contract--E reset
were executed.  The registry mismatch found during the audit has been repaired:
the registry now names the exported
\texttt{canonical\_value\_and\_analytical\_score}, and the governance test
resolves every registered entry point to a callable.  This closes discovery
only; it does not implement KDM or prove numerical score validity.  The existing
Section~3.6 prior-weight sidecar is diagnostic work from an isolated worktree;
it is not an integrated route in this checkout and is not a Younis IWSG
implementation through the reset.  The analytical function defaults to
\texttt{reset\_policy=none}, a diagnostic slice.  A canonical baseline
invocation must explicitly select \texttt{reset\_policy=contract\_e} and the
owner-registered dual-cap and trust-region settings.

\section{What the Younis construction supplies}

Younis and Sudderth begin with a weighted empirical cloud and replace its
Dirac representation by a regularized mixture
\begin{equation}
 \mu_h^N(\dd z)
   = \sum_{i=1}^{N} w_i k_{B_i}(z-x_i)\,\dd z,
 \qquad
 m_\theta(z)=\sum_{i=1}^{N}
       w_i(\theta)k_{B_i(\theta)}(z-x_i(\theta)).
 \label{eq:kdm}
\end{equation}
For their mixture-density particle filter, this mixture replaces discrete
resampling.  With a common Gaussian bandwidth one may generate an anchor
sample by
\begin{equation}
 I_j\sim\operatorname{Cat}(w_1,\ldots,w_N),\qquad
 z_j=x_{I_j}+L\varepsilon_j,\qquad
 \varepsilon_j\sim N(0,I),\quad LL^\mathsf{T}=B.
 \label{eq:younis-sampling}
\end{equation}
The 2024 construction stratifies the uniforms used for the categorical draw.
The label $I_j$ generates the sample, but the differentiable density is the
marginal sum in \eqref{eq:kdm}; it is not the density of the selected
component.  Thus every component can contribute to the gradient at every
sample location.  This is the source of the $N$-by-$N$ work in a
gradient-bearing resampling step.

Their IWSG construction fixes a proposal at the current parameter value.  With
\begin{equation}
 q_0(z)=m_{\theta_0}(z),\qquad
 \rho_j(\theta)=\frac{m_\theta(z_j)}{q_0(z_j)},\qquad
 z_j\sim q_0,
 \label{eq:iwsg-ratio}
\end{equation}
the finite estimator of a mixture expectation and its derivative are
\begin{align}
 \widehat I(\theta)
   &=\frac{1}{M}\sum_{j=1}^{M}\rho_j(\theta)\phi_\theta(z_j),\\
 \dd\widehat I
   &=\frac{1}{M}\sum_{j=1}^{M}
       \left\{\dd\rho_j\,\phi_\theta(z_j)
              +\rho_j\,\dd\phi_\theta(z_j)\right\}.
 \label{eq:iwsg}
\end{align}
At $\theta=\theta_0$, $\rho_j=1$.  The locations $z_j$ are held fixed in
this differentiation; parameter dependence is represented by the mixture
density ratio.  This is not the total JVP of an inverse-CDF draw or of a
parameter-dependent component index.  Equations~\eqref{eq:iwsg-ratio}--
\eqref{eq:iwsg} establish the derivative of a mixture expectation.  They do
not by themselves establish the derivative of an arbitrary nonlinear
operation that couples the complete set $(z_1,\ldots,z_M)$.  For such an
operation $H_\theta$, ordinary importance sampling on the joint sample would
instead use
\begin{equation}
 R_M(\theta)=\prod_{j=1}^M\rho_j(\theta),\qquad
 \dd\{R_M H_\theta\}\big|_{\theta_0}
 =\dd H_{\theta_0}+H_{\theta_0}
   \sum_{j=1}^M\dd\log m_\theta(z_j)\big|_{\theta_0}.
 \label{eq:joint-mixture-ratio}
\end{equation}
This joint ratio is exact for an expectation under independent mixture draws,
but its variance can grow rapidly with $M$.  Feeding individual IWSG weights
through a nonlinear \OT{} reset defines a finite weighted-particle algorithm;
it is not silently promoted to the joint identity
\eqref{eq:joint-mixture-ratio}.

The ratio in \eqref{eq:iwsg-ratio} is a raw importance weight.  Replacing it
by the self-normalized ratio
\begin{equation}
 \bar\rho_j(\theta)=\frac{\rho_j(\theta)}{\sum_l\rho_l(\theta)}
 \label{eq:self-normalized-ratio}
\end{equation}
changes the finite estimator.  At the anchor,
$\dd\log\rho_j=\dd\log m_\theta(z_j)$, whereas
\begin{equation}
 \dd\log\bar\rho_j
 =\dd\log m_\theta(z_j)
  -\frac1M\sum_l\dd\log m_\theta(z_l).
 \label{eq:centered-ratio-tangent}
\end{equation}
The second expression has been centered by a random sample average and is not
Equation~(15) of Younis and Sudderth.  Their released MDPF code confirms the
same ordering: it constructs the unnormalized gradient-injection factor
$\exp\{\log m-\operatorname{stop}(\log m)\}$, multiplies that factor by the
next measurement and dynamics weights, and only then normalizes the posterior
weights \citep{youniscode2023}.  The source code's small additive numerical
constant is an engineering detail and is not part of the mathematical route
defined here.

For a Gaussian component, let $u_i=z-x_i$.  Its complete directional
differential is
\begin{equation}
 \dd\log k_{B_i}(u_i)
 =
 u_i^\mathsf{T}B_i^{-1}\dd x_i
 +\frac12\left[
   u_i^\mathsf{T}B_i^{-1}(\dd B_i)B_i^{-1}u_i
   -\operatorname{tr}(B_i^{-1}\dd B_i)
 \right].
 \label{eq:component-tangent}
\end{equation}
Therefore
\begin{equation}
 \dd m_\theta(z)
 =\sum_i\left[
       (\dd w_i)k_{B_i}(u_i)
       +w_i k_{B_i}(u_i)\dd\log k_{B_i}(u_i)
     \right]. \label{eq:mixture-tangent}
\end{equation}
Equation~\eqref{eq:mixture-tangent} is the term that a KDM score path must
carry for every cloud, weight, or bandwidth dependence declared to vary with
$\theta$.  A bandwidth or cloud may instead be a genuinely fixed numerical
design, in which case its zero tangent is part of the scalar's definition.
Stopping a dependence that the declared scalar retains is a partial
derivative; declaring a quantity fixed before defining the scalar is not.

The source papers establish this mixture-gradient construction for learned,
discriminative state-density objectives.  In particular, the 2023 MDPF uses
the mixture both for posterior-state loss evaluation and for resampling; the
2024 paper again gives the marginalized ratio and separates resampling and
posterior bandwidths.  Neither paper derives the observed-data score of a
generative state-space model, an \LEDH{} flow, an \OT{} reset, \GenUT{}
moments, or dual-cap constraints.  Their complexity discussion makes the
computational boundary explicit: evaluating the mixture gradient at every
resampled location creates all-pairs interactions.  The papers' linear-cost
inference path sets forward resampling weights to uniform and does not compute
these gradients.  It cannot be transplanted to a score path without a new
approximation and error analysis.

\section{Why positive bandwidth is a target change}

The atom cloud is
\begin{equation}
 \mu_0^N(\dd z)=\sum_i w_i\delta_{x_i}(\dd z).
 \label{eq:atom-measure}
\end{equation}
For a test function $\phi$,
\begin{equation}
 \int\phi(z)\mu_h^N(\dd z)
 =\sum_i w_i\int\phi(x_i+u)K_{B_i}(\dd u),
 \label{eq:smoothed-functional}
\end{equation}
which is not generally $\sum_i w_i\phi(x_i)$.  This gives the following
simple boundary.

\begin{proposition}
If at least one component kernel has nonzero bandwidth, then, in general,
the KDM functional and the atom functional are different.  Consequently a
positive-bandwidth derivative cannot be called $S_0^N$ without an exact
correction or a separately proved limiting argument.
\end{proposition}

\begin{proof}
If the two functionals were equal for every smooth test function, convolution
by every component kernel would act as the identity on the values at its
component mean.  A non-degenerate Gaussian (or any ordinary smoothing
kernel) does not have this property.  Choose a smooth function with nonzero
curvature in a direction with positive kernel variance; its kernel average
differs from its value at the mean.  The finite sums then differ unless an
accidental cancellation occurs.  Differentiating the two unequal functionals
cannot restore equality for all parameter directions.
\end{proof}

The case $B=0$ is an explicit atom branch, not a request to evaluate a
singular Gaussian density.  A zero-bandwidth parity test must therefore not
form a Cholesky factor or an inverse of $B=0$ and then infer that the result is
the atom program.

For a linear-Gaussian observation model, a Gaussian state kernel gives the
useful exact convolution
\begin{equation}
 \overline g_\theta(y\mid x,B)
 = \int g_\theta(y\mid x+u)N(u;0,B)\,\dd u
 = N\!\left(y;C_\theta x,\,
       R_\theta+C_\theta B C_\theta^\mathsf{T}\right).
 \label{eq:linear-convolution}
\end{equation}
Write
\begin{align}
 S_i &= R_\theta+C_\theta B_i C_\theta^\mathsf{T}, &
 \mu_i &= C_\theta x_i, &
 r_i &= S_i^{-1}(y-\mu_i).
\end{align}
For fixed data $y$, the total directional differential used below is
\begin{align}
 \dd\mu_i
   &= (\dd C)x_i+C\,\dd x_i,\\
 \dd S_i
   &= \dd R+(\dd C)B_iC^\mathsf{T}
      +C(\dd B_i)C^\mathsf{T}+CB_i(\dd C)^\mathsf{T},\\
 \dd\log \overline g_i
   &=r_i^\mathsf{T}\dd\mu_i
     +\frac12\left\{
       r_i^\mathsf{T}(\dd S_i)r_i
       -\operatorname{tr}(S_i^{-1}\dd S_i)
     \right\}.
 \label{eq:linear-convolution-tangent}
\end{align}
Thus state, observation-map, observation-covariance, and bandwidth tangents
are all required when those quantities vary.  Positive semidefinite $B_i$ is
permitted here, including rank-deficient $B_i$, because only the strictly
positive definite observation covariance $S_i$ is factored.  This statement
does not assert an ambient Lebesgue density for a singular state kernel.
The corresponding observation increment
\begin{equation}
 \overline Z_t
 =\sum_i w_{t,i}^{-}\overline g_\theta(y_t\mid x_{t,i}^{-},B_{t,i}),
 \qquad
 \overline L_h^N=\sum_t\log\overline Z_t
 \label{eq:smoothed-value}
\end{equation}
is a well-defined new finite program.  It equals the atom increment only at
zero bandwidth or in special accidental cases.  If it is fed into the
subsequent reset, it defines an even larger new program.

The same distinction applies to a KDM placed after the reset.  Replacing
$x_t^+$ by a draw from \eqref{eq:kdm} changes the input to the next
transition and therefore changes all later normalizers.  The change may be
scientifically useful, but it is not an implementation detail.

\section{Legitimate ways KDM can enter}

\subsection{An auxiliary estimator or control variate}

The lowest-risk use keeps \eqref{eq:atom-program} and its value unchanged.
A KDM is formed from the actual cloud and produces a separately named mixture
score or force.  With the canonical finite value retained as the energy, the
mixture field is a proposal field, not the canonical score.  A control-variate
version must retain the exact residual:
\begin{equation}
 L_0^N=A_h+\left(L_0^N-A_h\right),\qquad
 \dd L_0^N=\dd A_h+\dd\left(L_0^N-A_h\right).
 \label{eq:residual}
\end{equation}
Dropping the residual changes the target.  A coefficient estimated on
calibration paths must be frozen before validation, and its uncertainty must
be measured on disjoint paths.  This route tests whether the KDM contains
useful score information without silently changing the filter.

\subsection{A KDM proposal with an exact model-target ratio}

There is one direct KDM use that leaves the generative model factors in the
importance ratio, but it is a BayesFilter extension rather than the Younis
MDPF.  The treatment of ancestry must be fixed before writing the ratio.  If
an ancestor $A$ is retained, draw it with a frozen probability $r_{t,0}(a)$,
then draw a pre-flow state from a complete conditional mixture
$q^{\KDM}_{t,0}(x_0\mid a,y_t)$.  Apply the ancestor-specific LEDH map
$x_1=F_{t,\theta,a}(x_0)$, with forward Jacobian $J_{t,a}$.  The induced joint
proposal density is
\begin{equation}
 q_t(j,x_1)
 = r_{t,0}(j)q^{\KDM}_{t,0}(x_0\mid j,y_t)
   \,|\det J_{t,j}(x_0)|^{-1},
 \qquad x_0=F_{t,\theta,j}^{-1}(x_1).
 \label{eq:kdm-proposal-density}
\end{equation}
The corresponding unnormalised weight is
\begin{equation}
 \widetilde w_t
 = \frac{w^-_{t-1,J}
     f_\theta(x_1\mid x^+_{t-1,J})
     g_\theta(y_t\mid x_1)
     |\det J_{t,J}(x_0)|}
    {r_{t,0}(J)q^{\KDM}_{t,0}(x_0\mid J,y_t)}.
 \label{eq:kdm-proposal-weight}
\end{equation}
Here ``complete'' means the sum over every kernel component within the
conditional mixture for the retained ancestor, not the kernel of the component
that generated $x_0$.  If ancestry is marginalized instead, the numerator must
also be marginalized,
\begin{equation}
 \widetilde w_t(x_0)
 =\frac{g_\theta(y_t\mid x_1)|\det J_t(x_0)|
   \sum_a w^-_{t-1,a}f_\theta(x_1\mid x^+_{t-1,a})}
  {q^{\KDM}_{t,0}(x_0)}.
 \label{eq:kdm-proposal-marginal-weight}
\end{equation}
This second expression is valid only when the flow map does not depend on an
ancestor that has been marginalized.  Mixing a marginalized denominator with
the conditional numerator in \eqref{eq:kdm-proposal-weight} is wrong.

An auxiliary-particle lookahead may guide $r_{t,0}$, but its approximate
predictive likelihood cannot replace either the exact transition factor
$f_\theta$ or the observation factor $g_\theta$.  Under the usual support,
invertibility, and differentiation-under-the-integral conditions, the sample
average $\widehat Z_t$ and its derivative are unbiased for the one-step
particle predictive integral conditional on the incoming cloud.  The ratio
$\dd\widehat Z_t/\widehat Z_t$ is generally biased for its logarithmic
derivative.  Moreover, the deterministic Contract--E reset changes the next
incoming cloud, so this one-step identity is not a proof that the complete
finite-horizon algorithm is an unbiased likelihood or score estimator.

Following the fixed-denominator principle of Younis IWSG, this conditional
LEDH extension freezes $r_{t,0}$ and $q^{\KDM}_{t,0}$ at an anchor $\theta_0$ while
differentiating the model numerator and LEDH map.  Younis and Sudderth do not
derive this conditional LEDH construction.  If an implementation instead
recomputes the proposal as $\theta$ changes, it needs a separate derivation
that couples the sampling map, proposal density, ancestor selection, and state
map.  Merely subtracting proposal or selection log-density tangents from a
fixed-sample importance weight is wrong for the anchored model-normalizer
derivative.  This arm is therefore labelled
\texttt{MODEL-IS} and is compared separately with both the canonical
\texttt{ATOM-FINITE} score and the exact Kalman score.

At the atom reset itself, an ordinary density ratio cannot repair the change:
a finite atom measure and an absolutely continuous KDM are mutually singular.
The KDM proposal must be inserted before a continuous transition, or the route
must define a new measure and a new finite target explicitly.

\subsection{Phase 4A: integrated kernelized observation weighting}

The first integrated experiment replaces only the point observation factor in
\eqref{eq:prior-observation-factorization}.  With the ordinary
pre-observation PF--PF logit $b_{t,i}$, define
\begin{align}
 a^h_{t,i}&=b_{t,i}+\log\overline g_\theta(
   y_t\mid x_{t,i}^{\mathrm{out}},B_{t,i}),\\
 Z^h_t&=\sum_i\exp(a^h_{t,i}), &
 \alpha^h_{t,i}&=\exp(a^h_{t,i})/Z^h_t,\\
 L_{\mathrm{obsKDM}}^N&=\sum_t\log Z^h_t, &
 S_{\mathrm{obsKDM}}^N&=\dd L_{\mathrm{obsKDM}}^N.
 \label{eq:observation-kdm-program}
\end{align}
Its one-step tangent is
\begin{align}
 \dd\log Z^h_t
 =\sum_i\alpha^h_{t,i}
   \left(\dd b_{t,i}+\dd\log\overline g_{t,i}\right),\qquad
 \dd\alpha^h_{t,i}
 =\alpha^h_{t,i}\left[
   \dd a^h_{t,i}-\dd\log Z^h_t\right].
 \label{eq:observation-kdm-weight-tangent}
\end{align}
The executed program passes $(x_t^{\mathrm{out}},\dd x_t^{\mathrm{out}},
\alpha_t^h,\dd\alpha_t^h)$ into the existing Contract--E reset and the
existing \GenUT{}/dual-cap correction.  Their output and tangent feed the next
time step.  Equations~\eqref{eq:linear-convolution-tangent} and
\eqref{eq:observation-kdm-weight-tangent} therefore define a total derivative
of this complete finite program, not a conditional observation derivative.

The reset also carries the particle-local UKF covariance with the same
target-by-source transport used for the state.  If $A_{ri}$ is the normalized
Sinkhorn transport from source particle $i$ to reset particle $r$, then the
covariance recursion and its total tangent are
\begin{align}
 P'_{r} &= \sum_i A_{ri}P_i, &
 \dd P'_{r} &= \sum_i (\dd A_{ri})P_i+
                    \sum_i A_{ri}\dd P_i.\label{eq:covariance-carry}
\end{align}
The carried $P'_r$, not the old source ordering, is consumed by the next UKF
prediction.  Here $P_i$ denotes a particle-local conditional covariance; no
between-cloud scatter term is added.  This same-transport carry is the
declared Contract--E extension of categorical triple resampling.  A previous
implementation omitted it at the analytical endpoint; that call-chain bug
was repaired and is now covered by endpoint and two-step tests.

This program is deliberately called \emph{kernelized observation weighting}.
It is not the complete Bayesian update of a Gaussian mixture: it does not
propagate the conditional posterior mean and covariance of every mixture
component.  It is consequently not a full Younis mixture-density particle
filter.  Positive bandwidth is labelled \texttt{KDM-FINITE}; the explicit
$B=0,\dd B=0$ branch is labelled \texttt{ATOM-FINITE} and must reproduce the
canonical trajectory.

A second insertion point represents the post-reset cloud by a continuous
mixture before the next transition.  It defines the distinct
\texttt{RESKDM-IWSG-FINITE} scalar and is derived and implemented separately
in the next subsection; it is not part of the Phase~4A
observation-weighting scalar.

The first paired experiment now answers the narrow Phase~4A question for one
small linear--Gaussian cell.  With $N=32$, $T=5$, 20 calibration paths, and
100 disjoint validation paths, calibration selected $\rho=0.8$ in
$B=\rho^2Q$.  The atom score MSE was $3.3366$ and the kernelized-observation
MSE was $3.4980$, a relative change of $-4.84\%$ when positive values denote
improvement.  The paired 95\% bootstrap interval for the MSE difference was
$[-0.223,0.597]$, so the two scores are not statistically ranked.  Inspection
of all prespecified bandwidths on this now-consumed holdout showed no positive
point improvement; $\rho=1.6$ was worse with interval $[0.189,2.399]$.

The mechanism is visible in the score errors.  From $\rho=0$ to $\rho=0.8$,
their sample standard deviation decreased from $1.616$ to $1.580$, but their
mean moved from $-0.868$ to $-1.014$.  The modest variance reduction was more
than offset by additional bias.  This rejects Phase~4A for promotion in this
scope.  It does not test the complete Younis resampling operation described
next, and it does not license tuning on the inspected validation paths.

For a finite IWSG normalizer,
\begin{align}
 \widehat Z_h(\theta)
   &=\frac{1}{M}\sum_j
       \rho_j(\theta)g_\theta(y_t\mid z_j),\\
 \dd\log\widehat Z_h
   &=\frac{\sum_j\rho_jg_j
       \left(\dd\log\rho_j+\dd\log g_j\right)}
           {\sum_j\rho_jg_j}.
 \label{eq:iwsg-normalizer}
\end{align}
Equation~\eqref{eq:iwsg-normalizer} is a derivative of the KDM
importance-sampling program.  It is not an identity with the atom program.

\subsection{Phase 4B: a complete mixture-resampling reference}

The source operation can now be stated without an ancestry ambiguity.  Let
$(c_{t,i},\alpha_{t,i},P_{t,i})$ denote the locations, normalized weights, and
particle-local UKF covariance marks presented to a continuous resampler.  If
the KDM replaces Contract--E, these are the weighted post-flow particles.  In
the combined BayesFilter reference selected below, they are the Contract--E
output, its carried covariances, and uniform weights
$\alpha_{t,i}=1/N$ with $\dd\alpha_{t,i}=0$.  In either case define the
marginalized resampling density
\begin{equation}
 m_{t,\theta}(z)=\sum_{i=1}^N
   \alpha_{t,i}(\theta)k_{B_t(\theta)}
      (z-c_{t,i}(\theta)).
 \label{eq:phase4b-mixture}
\end{equation}
The first reference uses the source paper's common Gaussian bandwidth and an
ordinary full-rank linear--Gaussian fixture.  The complete sequential route is
not called source-faithful because its Contract--E placement and covariance
marks are BayesFilter extensions.  A DSGE version must replace this density by
the chart density in Section~6.

At an anchor $\theta_0$, stratified categorical uniforms and Gaussian noises
generate fixed locations $z_{t,j}$ from $m_{t,\theta_0}$.  Component labels are
recorded to replay the draw, but are marginalized in every density evaluation.
For a nearby $\theta$, define
\begin{align}
 r_{t,j}(\theta)
   &=\frac{m_{t,\theta}(z_{t,j})}
           {m_{t,\theta_0}(z_{t,j})}, &
 \eta_{t,j}(\theta)
   &=\frac{1}{N}r_{t,j}(\theta),
 \label{eq:phase4b-iwsg-weights}\\
 \gamma_{t,ji}(\theta)
   &=\frac{\alpha_{t,i}k_{B_t}(z_{t,j}-c_{t,i})}
           {m_{t,\theta}(z_{t,j})}. &&
 \label{eq:phase4b-responsibilities}
\end{align}
At the anchor, $r_{t,j}=1$ and $\eta_{t,j}=1/N$.  The locations are fixed in the IWSG
differential, so $\dd z_{t,j}=0$, while
\begin{align}
 \dd\log r_{t,j}
   &=\sum_i\gamma_{t,ji}
      \left\{\dd\log\alpha_{t,i}
             +\dd\log k_{B_t}(z_{t,j}-c_{t,i})\right\},\\
 \dd\log\eta_{t,j}&=\dd\log r_{t,j}, &
 \dd\eta_{t,j}&=\eta_{t,j}\dd\log r_{t,j}.
 \label{eq:phase4b-weight-tangent}
\end{align}
Equations~\eqref{eq:component-tangent} and
\eqref{eq:phase4b-weight-tangent} include the cloud, weight, and bandwidth
terms.  Evaluating all $\gamma_{t,ji}$ is the required $N^2$ operation.  The
sum of the $\eta_{t,j}$ tangents is generally nonzero for a finite sample;
forcing it to zero would replace the raw IWSG correction by
\eqref{eq:self-normalized-ratio}.

Younis's particles carry locations and weights, not the UKF covariance used by
the LEDH implementation.  The following mark transport is therefore an
explicit BayesFilter extension:
\begin{align}
 \widetilde P_{t,j}
   &=\sum_i\gamma_{t,ji}P_{t,i}, &
 \dd\widetilde P_{t,j}
   &=\sum_i(\dd\gamma_{t,ji})P_{t,i}
     +\sum_i\gamma_{t,ji}\dd P_{t,i},\\
 \dd\gamma_{t,ji}
   &=\gamma_{t,ji}\left\{
       \dd\log\alpha_{t,i}+\dd\log k_{t,ji}
       -\dd\log m_{t,\theta}(z_{t,j})\right\}.
 \label{eq:phase4b-covariance-mark}
\end{align}
No between-component scatter is added: the kernel perturbation has already
been realized in $z_{t,j}$, and $P_{t,i}$ remains the local UKF algorithmic
mark.  This choice is testable, but it is not a result from the source paper.
The prespecified comparator instead carries the covariance mark of the fixed
anchor label,
\begin{equation}
 \widetilde P_{t,j}^{\mathrm{label}}=P_{t,I_{t,j}},\qquad
 \dd\widetilde P_{t,j}^{\mathrm{label}}=\dd P_{t,I_{t,j}}.
 \label{eq:phase4b-selected-mark}
\end{equation}
The label has no tangent.  Both policies retain the same marginalized mixture
density and raw IWSG weights; only the downstream UKF mark differs.  Neither
mark policy is supplied by Younis and Sudderth, so their comparison diagnoses a
BayesFilter extension rather than source fidelity.

There are now two possible placements, and they must not share a name.
Replacing Contract--E by mixture resampling gives the source MDPF placement
but removes the requested Contract--E, \GenUT{}, and dual-cap reset.  Retaining
Contract--E and then applying \eqref{eq:phase4b-iwsg-weights} before the next
transition executes both mechanisms.  The latter is the declared Phase~4B
BayesFilter reference.  Its fixed locations, raw IWSG weights, and covariance
marks become the parents of the next transition; all later PF--PF,
flow-Jacobian, reset, and cap terms are then differentiated normally.

To make the finite scalar explicit, let $h_{t+1,j}(\theta)$ contain the full
next-step PF--PF multiplier for child $j$: transition density, observation
density, flow Jacobian, and proposal denominator.  The post-resampling
increment is
\begin{equation}
 Z_{t+1}^{\mathrm{IWSG}}(\theta)
 =\sum_{j=1}^N\eta_{t,j}(\theta)h_{t+1,j}(\theta)
 =\frac1N\sum_{j=1}^N r_{t,j}(\theta)h_{t+1,j}(\theta).
 \label{eq:phase4b-next-normalizer}
\end{equation}
The posterior weights are normalized only after the product
$\eta_{t,j}h_{t+1,j}$ is formed.  This preserves the raw importance-sampling
normalizer in \eqref{eq:phase4b-next-normalizer}.  Call the resulting scalar
\begin{equation}
 L_{\mathrm{resKDM,IWSG}}^N(\theta;\xi,z^0,q^0)
 =\sum_{t=1}^T\log Z_t^{\mathrm{IWSG}},\qquad
 S_{\mathrm{resKDM,IWSG}}^N=\dd L_{\mathrm{resKDM,IWSG}}^N.
 \label{eq:phase4b-finite-target}
\end{equation}
This is the combined \texttt{RESKDM-IWSG-FINITE} program.  It is neither
$L_0^N$ nor the Younis MDPF alone.  A two-pass evaluation is required: an
anchor pass generates every sequential $z_{t,j}$ and denominator
$m_{t,\theta_0}(z_{t,j})$; a replay pass evaluates the same fixed sample bank
and its total tangent for other $\theta$.  Re-anchoring inside an HMC
trajectory would change the finite energy and is forbidden.

The raw IWSG weights are part of this new finite scalar, not auxiliary
metadata.  In particular, if $\omega_{t+1,j}$ denotes the incoming weight used
by the next PF--PF step, then
\begin{equation}
 \log\omega_{t+1,j}=\log r_{t,j}-\log N,\qquad
 \dd\log\omega_{t+1,j}=\dd\log r_{t,j}.
 \label{eq:phase4b-incoming-weight}
\end{equation}
The next normalizer must consume both quantities.  Resetting the numerical
weights to $1/N$ while retaining only the sampled locations, or retaining the
weight values while discarding their tangent, removes the IWSG mechanism.

Prematurely normalizing the ratios defines a different scalar.  Because the
subsequent posterior softmax removes common scale, the state and posterior
weight recurrences agree, but the accumulated normalizer does not.  With the
same samples and covariance-mark policy,
\begin{equation}
 L_{\mathrm{resKDM,SN}}^N
 =L_{\mathrm{resKDM,IWSG}}^N
  -\sum_{t=1}^{T-1}\log\left\{\frac1N\sum_{j=1}^N r_{t,j}\right\}.
 \label{eq:phase4b-self-normalized-difference}
\end{equation}
The first Phase~4B implementation computed this
\texttt{RESKDM-SN-FINITE} scalar.  Its finite-difference agreement was correct
for that scalar but did not validate the intended raw-IWSG score.  The repaired
route implements \eqref{eq:phase4b-incoming-weight}; an executable regression
requires a generally nonzero sum of raw weight tangents and rejects the
centered tangent in \eqref{eq:centered-ratio-tangent} as its incoming score.

There are also two distinct meanings of ``unbiased'' that must not be
conflated.  Equations~\eqref{eq:iwsg-ratio}--\eqref{eq:iwsg} give the ordinary
importance-sampling identity for a single mixture expectation.  For a
nonlinear function of the complete resampled cloud, the exact joint identity
is \eqref{eq:joint-mixture-ratio}.  The sequential use of raw individual
weights in \eqref{eq:phase4b-incoming-weight} follows the IWSG
gradient-injection ordering, while using their scale in a generative evidence
normalizer is a BayesFilter extension.  Its analytical derivative is the total
derivative of $L_{\mathrm{resKDM,IWSG}}^N$ as defined here, but neither the
source paper nor the present derivation proves that it is an unbiased estimator of
$\nabla_\theta\log p_\theta(y_{1:T})$, of the expected stochastic-filter loss,
or of $S_0^N$.  Those are empirical or additional-theory questions.

The zero-bandwidth limit requires separating the locations from their IWSG
derivative.  In this combined route, the component weights are exactly $1/N$
and the $j$th stratified position is $(j-1+u_j)/N$, with $0\le u_j<1$.
Consequently the selected label is $I_j=j$: each component is selected once.
The earlier statement that this particular route tends to repeated component
locations was wrong.  With $B=h^2Q$ and $L_QL_Q^\mathsf{T}=Q$, its anchor
sample is $z_j=c_j+hL_Q\varepsilon_j$, which tends to $c_j$.

This location limit does not establish score parity.  For distinct centres,
fixed positive-definite $Q$, and a parameter-independent bandwidth scale $h$,
the selected component dominates the mixture as $h\to0$.  At the anchor,
its leading location contribution to the IWSG differential is
\begin{equation}
 \dd\log m_\theta(z_j)\big|_{\theta_0}
 \simeq \frac{1}{h}\varepsilon_j^\mathsf{T}L_Q^{-1}\dd c_j.
 \label{eq:phase4b-small-bandwidth-tangent}
\end{equation}
Thus the fixed-sample tangent can grow while the forward locations converge.
Repeated centres require a separate mixture-limit analysis.  At $B=0$ the
ambient Gaussian density is undefined, so Phase~4B does not implement a
zero-bandwidth score branch equal to \texttt{ATOM-FINITE}.  Its checks are instead anchor
identity, complete-mixture density, analytical-versus-finite-difference
tangents, sequential replay, and comparison with both the canonical score and
the Kalman score.

The exact all-components reference costs $O(N^2)$ per mixture-gradient
evaluation and $O(TN^2)$ for a sequential gradient.  Sparse neighbors, random
features, and low-rank approximations would be new algorithms; they are not
permitted in the reference and require separate error analyses if considered
later.

\section{Degenerate DSGE support}

A full-dimensional kernel is inappropriate when the DSGE transition has
deterministic directions.  A first constrained construction injects noise only
through the stochastic innovation map:
\begin{equation}
 u=G_\theta L_\beta\varepsilon,\qquad
 B_\theta=G_\theta H_\beta G_\theta^\mathsf{T},\qquad
 \varepsilon\sim N(0,I).
 \label{eq:innovation-kernel}
\end{equation}
For a linear support constraint $D x=c$, the kernel must satisfy
$D G_\theta=0$.  For a nonlinear manifold, a chart or retraction
$x=\rho_\theta(v,\varepsilon)$ is required, together with the appropriate
coordinate-space density and Jacobian.  When $B_\theta$ is rank deficient,
there is no ordinary full-dimensional Lebesgue density.  The implementation
must evaluate the mixture in innovation or chart coordinates rather than
pretending that $B_\theta$ is a full SPD ambient covariance.

This constrained kernel is an extension motivated by the DSGE target, not a
claim from the Younis papers.  Its first tests are rank, constraint residual,
and finite-difference tests for the $G_\theta$ and $H_\beta$ tangents.  A
positive bandwidth that leaves the support is a hard failure even if its
numerical score looks smoother.

The Phase~2 diagnostic implements the fixed-chart special case: a shared
orthonormal basis $U\in\R^{d\times r}$ is frozen, the Gaussian mixture is
evaluated in $U^\mathsf{T}x$ coordinates, and the orthogonal residual is checked
for both values and tangents.  This gives a rank-aware density with respect to
the chart measure for a fixed linear support.  It does not yet differentiate a
parameter-dependent $G_\theta$, chart, or retraction; those Jacobian terms are
required before a DSGE production route can be considered.

\section{Bounded validation program}

The first diagnostic is a TensorFlow reconstruction of the Younis mixture
algebra with an analytical tangent.  The proposal-corrected LEDH arm is a
separate extension, not a source theorem.  The evaluator must expose cloud,
weight, and bandwidth derivatives and provide a small all-pairs reference.  The
claim-bearing call chain must resolve from the actual canonical endpoint to
this evaluator; a standalone function is not evidence of integration.

The following gates are required:
\begin{enumerate}[leftmargin=2em]
  \item For Phase~4A, at zero bandwidth, value and score match the canonical
        fixed-stream program within the declared numerical tolerance.
  \item The Phase~4A zero-bandwidth branch is an atom calculation; no Cholesky
        factor or inverse of a zero covariance is used to manufacture the
        limit.  Phase~4B makes no zero-bandwidth parity claim.
  \item At positive bandwidth, the result records
        $\overline L_h^N-L_0^N$ and is not labelled value-preserving.
  \item Finite differences of the declared $L_h^N$ agree with its complete
        analytical tangent when state, observation matrix, process and
        observation covariances, bandwidth, and model parameters are
        perturbed separately and jointly.
  \item A degenerate DSGE fixture has zero support violation and a correct
        rank-aware density calculation.
  \item Any route claiming linear cost has an explicit no-$N^2$ memory and
        work check.  The all-pairs implementation remains diagnostic.
  \item On a linear-Gaussian model, score error is compared with the Kalman
        score over paired paths and independent fixed streams.  Bias and
        variance are reported separately.
  \item For a \texttt{MODEL-IS} arm, the proposal-density identity is checked
        separately from the post-reset finite value; the complete mixture
        denominator, support, and forward Jacobian must all be present.
  \item For Phase~4B, every anchor sample is replayed exactly, its forward
        importance ratio is one, every component contributes to the
        marginalized density, and the $N^2$ responsibility and covariance-mark
        tangents agree with finite differences.
  \item The sequential Phase~4B endpoint is tested at $T\ge2$ so that the
        post-reset KDM state, weight, and covariance marks demonstrably affect
        a later normalizer.  A one-step density helper is not integration
        evidence.
  \item The next PF--PF step consumes the raw IWSG log weights and their
        uncentered tangents from \eqref{eq:phase4b-incoming-weight}; executable
        ablations must detect either one being discarded or prematurely
        normalized.
  \item The score-quality fixture has at least two state dimensions.  Its
        executed trace must contain a nonzero off-diagonal pairwise target mask,
        a finite nonzero pairwise correction displacement, and both pairwise and
        coordinate-cap diagnostics.  Merely setting nonzero options in a scalar
        model does not execute the pairwise algorithm.
\end{enumerate}

The first score-quality campaign freezes a two-state, transition-first model.
Writing $x_t=A(\theta)x_{t-1}+\epsilon_t$ and
$y_t=Hx_t+\nu_t$, its matrices are
\begin{align}
 A(\theta)
 &=
 \begin{bmatrix}0&0.12\\-0.08&0.10\end{bmatrix}
 +\theta\begin{bmatrix}1&0\\0&0.7\end{bmatrix},
 &\theta_0&=0.72,\\
 H&=\begin{bmatrix}1&0.25\\-0.15&0.9\end{bmatrix},
 &Q&=\begin{bmatrix}0.12&0.025\\0.025&0.09\end{bmatrix},\\
 R&=\begin{bmatrix}0.22&0.035\\0.035&0.18\end{bmatrix},
 &P_0&=\begin{bmatrix}0.8&0.12\\0.12&0.6\end{bmatrix}.
 \label{eq:phase4b-matrix-lgssm}
\end{align}
Here $x_{-1}\sim N(0,P_0)$ and the first observation follows one transition.
At $\theta_0$, $A(\theta_0)$ is stable.  An independent analytical Kalman
recursion differentiates its prediction, innovation covariance, gain, posterior
mean, and posterior covariance with $\dd A=\operatorname{diag}(1,0.7)$.  The
first particle scopes are $(N,T)=(32,5)$ and $(64,20)$.  The smaller scope runs
first, and the larger runs only if the scientific vetoes and total compute
budget permit it.  This matrix fixture replaces the earlier scalar proposal,
which was wrong for testing the full pairwise correction because that
correction is identically absent in one dimension.

The comparator ladder is:
\begin{enumerate}[leftmargin=2em]
  \item the current canonical analytical total JVP with no KDM;
  \item a constrained innovation-jitter proposal with no KDM score
        replacement;
  \item a plain bootstrap particle score on the non-degenerate
        linear-Gaussian reference;
  \item the KDM auxiliary route with the canonical score retained;
  \item the Phase~4A observation-KDM route with its new target; and
  \item the Phase~4B post-reset IWSG-resampling route with its new target.
\end{enumerate}
The Kalman solution is an oracle, not a heuristic.  Comparisons must be
conditional on short and long horizons, degenerate and non-degenerate
support, separated and overlapping modes, and small and large particle
counts.  Effective sample size, mode coverage, runtime, and bandwidth
sensitivity explain a result; they do not establish score correctness.
The completed Phase~4A small-cell diagnostic contains only the exact Kalman
oracle and the canonical atom comparator.  It cannot satisfy this full
promotion ladder and is reported only as a rejection in its tested scope.

\section{Code boundary and next step}

The current canonical endpoints are the two functions named in
Section~2, with the source-route initialization in
\texttt{bayesfilter/highdim/ledh\_pfpf\_genut\_initial\_rqmc\_tf.py}.  The
existing prior-weight sidecar captures a corrected pre-observation cloud, but
it does not implement the Younis IWSG mixture gradient through the integrated
\LEDH{}--\OT{}--\GenUT{} reset.  It must not be relabelled as that method.

The registry call-chain repair is complete.  Phase~1 is now implemented as the
diagnostic module
\path{bayesfilter/highdim/ledh_younis_kdm_tf.py}.  Its route ID is
\begin{center}
\texttt{ledh\_younis\_kdm\_algebra\_}\texttt{diagnostic\_v1}
\end{center}
Its global and
conditional mixture kernels expose the all-pairs density and complete cloud,
weight, and covariance tangents.  The source-matched fixed-proposal IWSG
kernel has no sample-location tangent and keeps the proposal density at its
supplied anchor.  The conditional kernel is labelled as a proposal-density
evaluator, not a model-score estimator.  The separate anchored
\texttt{MODEL-IS} PF--PF kernel contains every value term in
\eqref{eq:kdm-proposal-density}--\eqref{eq:kdm-proposal-weight} and requires
an explicit per-particle certificate that the forward map is invertible; a
finite log determinant alone is not accepted as that certificate.  Its
tangent excludes proposal and selection derivatives because their sampling
law is frozen at the anchor.  The module is diagnostic-only and is not in the
canonical registry or call chain.

The bounded Phase~1 test suite passes its complete-tangent, rank, anchor,
proposal, and XLA-smoke checks.  Phase~2A adds

\begin{itemize}[leftmargin=2em]
  \item the controlled one-step linear-Gaussian normalizer\\
        \texttt{make\_linear\_gaussian\_kdm\_}\\
        \texttt{normalizer\_kernel}, whose
        positive branch is \texttt{KDM-FINITE} and whose explicit zero branch
        is \texttt{ATOM-FINITE}; and
  \item an opt-in \texttt{return\_trace=True} surface on the analytical
        endpoint, used to build the sidecar from the actual post-flow cloud
        and weights under the Contract--E/dual-cap settings.
\end{itemize}

The Phase~2A finite-difference and endpoint tests pass, and the Phase~2B
fixed-chart support tests pass.  Phase~3 now provides
\texttt{canonical\_linear\_gaussian\_kdm\_auxiliary}.  It reads the actual
analytical trace, reconstructs the zero-bandwidth value and score from
\eqref{eq:prior-observation-factorization}, and reports a separately labelled
positive-bandwidth functional without feeding it into the reset.  In
particular, it uses $\bar w_t$, not $w_t^+$; using $w_t^+$ would apply the
current observation twice.

Phase~4A is now implemented in
\path{bayesfilter/highdim/ledh_younis_kdm_integrated_tf.py}.  The route
\texttt{ledh\_younis\_kdm\_integrated\_observation\_weighting\_v1} calls the
same internal finite-program executor as the registered canonical wrapper and
substitutes only the observation factor.  The code implements
\eqref{eq:linear-convolution-tangent}, feeds the resulting weights and
tangents through Contract--E and both correction caps, and returns the actual
multi-step trace.  Its fixed-shape factory defaults to XLA compilation.

The implementation audit also repaired three pre-existing total-derivative
gaps in the shared analytical executor: generic Gaussian transition and
observation fallbacks now include direct covariance-density tangents; a model
may supply the total tangent of its observation Jacobian; and every resulting
$\dd H$ term enters the LEDH flow.  The flow log determinant now uses an
XLA-compatible QR/triangular-solve identity rather than an unsupported matrix
determinant/inverse pair.  Rebuilt-model finite differences test these terms.

Phase~4A is a correct total derivative of its declared
\texttt{KDM-FINITE} scalar; it is wrong if described as the derivative of
\texttt{ATOM-FINITE}.  The current one-cell experiment gives no evidence that
it reduces model-score error, and the all-bandwidth inspection indicates that
its small variance reduction is offset by additional bias.  The complete
rows, manifest, and source hashes are recorded under
\path{docs/benchmarks/artifacts/ledh_younis_kdm_}\\
\path{phase4_20260908/campaign-attempt05-n32t5-all-rho-f32-current/}.

Phase~4B is no longer blocked by an unspecified mixture or ancestry rule:
Equations~\eqref{eq:phase4b-mixture}--\eqref{eq:phase4b-finite-target} define
the source-matched raw mixture ratio, the BayesFilter covariance marks,
the insertion point after Contract--E, and the two-pass anchor/replay rule.
The exact $N$-by-$N$ operation is implemented in
\path{bayesfilter/highdim/ledh_younis_kdm_resampling_tf.py}.  Its anchor and
replay endpoints call the shared canonical executor after Contract--E and both
correction caps, enforce one common full-rank Gaussian bandwidth at each time,
and expose both raw and self-normalized ratio diagnostics with complete
tangents.  The repaired route carries raw ratios divided by $N$, not
self-normalized ratios, into the next PF--PF normalizer.  The
responsibility-mean and fixed-label covariance-mark policies are separate,
explicit routes.  One-at-a-time and joint finite
differences, exact anchor replay, canonical no-hook parity, and isolated
two-step state/weight/covariance wiring tests pass.  Bounded float64 and
float32/no-TF32 GPU/XLA smokes of the earlier self-normalized route are
superseded for derivative certification.  On 9 September, four fresh smokes
of the repaired raw-IWSG program passed for both precisions and both covariance
mark policies at $D=2,N=8,T=2$.  These checks establish implementation consistency for
\texttt{RESKDM-IWSG-FINITE}; it supplies no evidence yet about
model-score error.  The separate conditional \texttt{MODEL-IS} proposal
remains an extension and is not a substitute for this combined program.
The canonical trace now retains, rather than discards, higher-moment validity,
the off-diagonal target mask, pairwise pre/post-cap displacement, pairwise cap
scale, and coordinate-cap activity.  The new matrix Kalman recursion and
fixed-stream bootstrap comparator pass analytical-versus-finite-difference
tests.  These are implementation checks; the powered campaign has not yet run.

Fresh Phase~4A endpoint calibration after the covariance-carry repair passes
for float64 and for float32 without TF32.  Its float32 TF32 arm still fails the
declared graph/XLA and atom-identity tolerances, so TF32 remains unadmitted for
that diagnostic route.  Phase~4B has so far been checked only in float64 and
float32 without TF32.  The reduced registered fused batch lane is also
excluded: its current body does not execute the
Contract--E/\GenUT{}/dual-cap reset.  No HMC promotion or canonical default
change is authorized.  A positive result at finite bandwidth would support a
new, explicitly named regularized program; it would not retroactively make its
derivative the score of \eqref{eq:atom-program}.

\bibliographystyle{plainnat}
\bibliography{ledh_younis_kdm_score}

\end{document}
